% GENERATED FILE. Edit canonical lessons, then run npm run generate:books.
% canonical-source-hash: fnv1a64:8364960e1e1628a3
% canonical-lessons: FR-C369-selon, FR-C369-eventuel, FR-C369-eventuellement, FR-C369-eventualite, FR-C369-faute-de

\chapter{Selon les cas}
\label{ch:fr-selon-les-cas}
\hlchaptermodality{369}

\begin{chapteropening}
\textbf{By the end of this chapter:} \emph{I can say that a plan depends on things: according to, a possible extra, possibly, an eventuality and for lack of.}

Complete the last lesson of chapter 369: I can say that a plan depends on things: according to, a possible extra, possibly, an eventuality and for lack of.
\end{chapteropening}

\section[selon]{selon — according to; depending on}

\label{lesson:FR-C369-selon}



\begin{quote}
Before the new one: say the French for in case, then the French for otherwise.
\end{quote}

\subsection*{You'll want to know: selon}
\textbf{selon} — \textquotedblleft{}according to; depending on\textquotedblright{}.

\begin{quote}
Selon le journal, il va neiger demain.
\end{quote}

\begin{quote}
\emph{According to the newspaper, it is going to snow tomorrow.}
\end{quote}

\emph{Selon moi} — in my view. \emph{C'est selon} — it depends.

\begin{cousinweb}
\emph{selon} comes from Old French \emph{selonc}, \emph{selunc} (EtymDB), usually traced to a Vulgar Latin \emph{sublongum} \textquotedblleft{}along\textquotedblright{}, built on \emph{longus} \textquotedblleft{}long\textquotedblright{} (Lewis \& Short).
\end{cousinweb}

\begin{grammarlens}[title={What you've built}]
One more word for telling, explaining and describing at B1.
\end{grammarlens}

\subsection*{Your turn}
\begin{itemize}
\raggedright
  \item \emph{Say it:} \emph{selon}
  \item \emph{Say it:} \emph{selon}, once more
  \item \emph{Say it:} a sentence of your own with \emph{selon}
  \item \emph{Recall:} say the French for in case, then the French for otherwise, then say \emph{selon} again
\end{itemize}

\begin{tcolorbox}[breakable,colback=teal!4,colframe=teal!35!black,title={Before you move on}]
What does selon mean? (\textquotedblleft{}According to; depending on\textquotedblright{}.) Say it once more.
\end{tcolorbox}
\section[éventuel]{éventuel — possible (if it happens)}

\label{lesson:FR-C369-eventuel}



\begin{quote}
Before the new one: say the French for otherwise, then the French for according to.
\end{quote}

\subsection*{You'll want to know: éventuel}
\textbf{éventuel} — \textquotedblleft{}possible (if it happens)\textquotedblright{}.

\begin{quote}
Nous avons gardé de l'argent pour les frais éventuels.
\end{quote}

\begin{quote}
\emph{We kept some money for any possible costs.}
\end{quote}

A false friend: \emph{éventuel} means possible, not final. Feminine \emph{éventuelle}.

\begin{cousinweb}
\emph{éventuel} is a learned formation on Latin \emph{ēventus} \textquotedblleft{}an outcome, an event\textquotedblright{} (Lewis \& Short), from \emph{ēvenīre} \textquotedblleft{}to come out, to happen\textquotedblright{}, built on \emph{venīre} \textquotedblleft{}to come\textquotedblright{}.
\end{cousinweb}

\begin{grammarlens}[title={What you've built}]
One more word for telling, explaining and describing at B1.
\end{grammarlens}

\subsection*{Your turn}
\begin{itemize}
\raggedright
  \item \emph{Say it:} \emph{éventuel}
  \item \emph{Say it:} \emph{éventuel}, once more
  \item \emph{Say it:} a sentence of your own with \emph{éventuel}
  \item \emph{Recall:} say the French for otherwise, then the French for according to, then say \emph{éventuel} again
\end{itemize}

\begin{tcolorbox}[breakable,colback=teal!4,colframe=teal!35!black,title={Before you move on}]
What does éventuel mean? (\textquotedblleft{}Possible (if it happens)\textquotedblright{}.) Say it once more.
\end{tcolorbox}
\section[éventuellement]{éventuellement — possibly, if need be}

\label{lesson:FR-C369-eventuellement}



\begin{quote}
Before the new one: say the French for according to, then the French for possible.
\end{quote}

\subsection*{You'll want to know: éventuellement}
\textbf{éventuellement} — \textquotedblleft{}possibly, if need be\textquotedblright{}.

\begin{quote}
Je peux éventuellement venir samedi.
\end{quote}

\begin{quote}
\emph{I can possibly come on Saturday.}
\end{quote}

A false friend: not eventually, which is \emph{finalement}.

\begin{cousinweb}
\emph{éventuellement} is built from \emph{éventuel} and \emph{-ment}. \emph{éventuel} rests on Latin \emph{ēventus} \textquotedblleft{}an outcome\textquotedblright{}, from \emph{ēvenīre} \textquotedblleft{}to happen\textquotedblright{} (Lewis \& Short).
\end{cousinweb}

\begin{grammarlens}[title={What you've built}]
One more word for telling, explaining and describing at B1.
\end{grammarlens}

\subsection*{Your turn}
\begin{itemize}
\raggedright
  \item \emph{Say it:} \emph{éventuellement}
  \item \emph{Say it:} \emph{éventuellement}, once more
  \item \emph{Say it:} a sentence of your own with \emph{éventuellement}
  \item \emph{Recall:} say the French for according to, then the French for possible, then say \emph{éventuellement} again
\end{itemize}

\begin{tcolorbox}[breakable,colback=teal!4,colframe=teal!35!black,title={Before you move on}]
What does éventuellement mean? (\textquotedblleft{}Possibly, if need be\textquotedblright{}.) Say it once more.
\end{tcolorbox}
\section[l'éventualité]{l'éventualité — the possibility, the eventuality}

\label{lesson:FR-C369-eventualite}



\begin{quote}
Before the new one: say the French for possible, then the French for possibly.
\end{quote}

\subsection*{You'll want to know: l'éventualité}
\textbf{l'éventualité} — \textquotedblleft{}the possibility, the eventuality\textquotedblright{}.

\begin{quote}
Il faut prévoir toutes les éventualités.
\end{quote}

\begin{quote}
\emph{We have to plan for every possibility.}
\end{quote}

\emph{Dans l'éventualité où} takes the conditional. Feminine.

\begin{cousinweb}
\emph{éventualité} is built from \emph{éventuel} and \emph{-ité}. \emph{éventuel} rests on Latin \emph{ēventus} \textquotedblleft{}an outcome\textquotedblright{}, from \emph{ēvenīre} \textquotedblleft{}to happen\textquotedblright{} (Lewis \& Short).
\end{cousinweb}

\begin{grammarlens}[title={What you've built}]
One more word for telling, explaining and describing at B1.
\end{grammarlens}

\subsection*{Your turn}
\begin{itemize}
\raggedright
  \item \emph{Say it:} \emph{l'éventualité}
  \item \emph{Say it:} \emph{l'éventualité}, once more
  \item \emph{Say it:} a sentence of your own with \emph{l'éventualité}
  \item \emph{Recall:} say the French for possible, then the French for possibly, then say \emph{l'éventualité} again
\end{itemize}

\begin{tcolorbox}[breakable,colback=teal!4,colframe=teal!35!black,title={Before you move on}]
What does l'éventualité mean? (\textquotedblleft{}The possibility, the eventuality\textquotedblright{}.) Say it once more.
\end{tcolorbox}
\section[faute de]{faute de — for lack of}

\label{lesson:FR-C369-faute-de}



\begin{quote}
Before the new one: say the French for possibly, then the French for the possibility.
\end{quote}

\subsection*{You'll want to know: faute de}
\textbf{faute de} — \textquotedblleft{}for lack of\textquotedblright{}.

\begin{quote}
Faute de temps, nous n'avons pas visité le musée.
\end{quote}

\begin{quote}
\emph{For lack of time, we did not visit the museum.}
\end{quote}

\emph{Faute de mieux} — for want of anything better.

\begin{cousinweb}
\emph{faute} is inherited through Old French from a Vulgar Latin \emph{fallita} \textquotedblleft{}a shortcoming\textquotedblright{} (EtymDB), from Latin \emph{fallere} \textquotedblleft{}to deceive, to fail\textquotedblright{} (Lewis \& Short).
\end{cousinweb}

\begin{grammarlens}[title={What you've built}]
One more word for telling, explaining and describing at B1.
\end{grammarlens}

\subsection*{Your turn}
\begin{itemize}
\raggedright
  \item \emph{Say it:} \emph{faute de}
  \item \emph{Say it:} \emph{faute de}, once more
  \item \emph{Say it:} a sentence of your own with \emph{faute de}
  \item \emph{Recall:} say the French for possibly, then the French for the possibility, then say \emph{faute de} again
\end{itemize}

\begin{tcolorbox}[breakable,colback=teal!4,colframe=teal!35!black,title={Before you move on}]
What does faute de mean? (\textquotedblleft{}For lack of\textquotedblright{}.) Say it once more.
\end{tcolorbox}
